\section{Conclusion}

\system turns fault-tolerant acquisition decisions into typed, auditable program structure.
The key idea is to keep the code index, linear resource context, conservative effect summary, backend side conditions, and certificate level visible while a logical program is elaborated into an executable acquisition plan.
In the running surface-code example, this visibility is what distinguishes a legal direct $\gate{H}$, an acquired $\gate{T}$, a hybrid switched region, and an illegal direct $\gate{CNOT}$.

The current implementation demonstrates the interface end to end across planning, checking, SMT obligations, protocol certificates, decoder contracts, geometry checks, and a Lean proof kernel.
The evidence is intentionally conditional: local algebraic and software-level obligations are checked locally, while complete physical protocol claims remain imported theorem dependencies.
This boundary is the paper's main programming-language message.
Certified local FTQC rules can be composed through a small, code-indexed calculus that makes legality, resources, effects, layout constraints, and proof assumptions explicit.
